% ================================================================
% RESULTS SECTION
% ================================================================

This section presents the estimation results in three stages. Section~\ref{sec:gravity_results} establishes that the gravity framework fits bilateral film production flows, with particular attention to how the determinants of film flows differ from goods trade. Section~\ref{sec:ppml_incentives} reports the PPML incentive estimates (Models~3--7), documenting the cross-sectional patterns in destination attraction and origin pull-back effects. Section~\ref{sec:pair_fe_results} subjects these findings to stricter identification via country-pair fixed effects (Models~8--11), revealing which effects survive and which reflect selection.

% ----------------------------------------------------------------
\subsection{The Gravity Framework for Film Production}
\label{sec:gravity_results}
% ----------------------------------------------------------------

Tables~\ref{tab:model1} and~\ref{tab:model2} report the baseline and full gravity specifications for both dependent variables. The gravity model fits bilateral film flows well, with pseudo-$R^2$ values of 0.62--0.66, though the determinants differ from goods trade in several notable respects.

Distance is negative and highly significant across all specifications, but the elasticity ($-0.24$ to $-0.28$ for film counts; $-0.18$ for budgets) is substantially smaller than the $-0.7$ to $-1.5$ range typically found in goods trade \citep{head2014}. This is consistent with the services trade literature, which finds that non-physical trade is less distance-sensitive \citep{walsh2006, kimura2006}. The smaller distance elasticity for budgets relative to film counts suggests that when distant productions do occur, they tend to be higher-budget---distance deters marginal low-budget projects more than large productions for which the right location justifies the cost.

Common language is the single most important predictor of bilateral film flows, increasing production by approximately 69\% for film counts and 47\% for budgets (Model~2). These effects far exceed the 15--30\% language effects typically found in goods trade and are consistent with the observation that film is an inherently language-intensive product: scripts, dialogue, direction, and crew communication all require linguistic compatibility. That common language matters more than distance is a distinctive feature of audiovisual services trade.

Contiguity increases film flows by 27\% for counts and 32\% for budgets. Two results are more surprising. Colonial relationships are associated with approximately 26\% \textit{less} bilateral film production, contrary to the positive colonial-tie effects found in goods trade. This may reflect that post-colonial cultural production has intentionally diverged. Regional trade agreements are associated with 17--21\% less production, possibly because RTA partners tend to be regional neighbours with independent domestic film industries rather than co-production complementarities.

Economic freedom is insignificant on both sides across all specifications. This is useful for identification: the incentive coefficients in subsequent models are not confounded by general institutional quality.

Origin GDP is consistently positive and significant, though modest for film counts (elasticity $\approx 0.05$) and much larger for budgets ($\approx 0.6$), reflecting that wealthier origin countries fund more and higher-value overseas productions. Destination GDP is insignificant for film counts and significantly negative for budgets, inverting the standard gravity prediction and suggesting that larger economies tend to be origins rather than destinations of production capital.

% ----------------------------------------------------------------
\subsection{Incentive Effects: PPML Estimates}
\label{sec:ppml_incentives}
% ----------------------------------------------------------------

Tables~\ref{tab:model3}--\ref{tab:model7} report the PPML incentive estimates. Adding incentive variables does not destabilise the gravity controls, which remain essentially unchanged from Model~2.

\subsubsection{Aggregate Incentive Effects}

Model~3 (Table~\ref{tab:model3}) introduces the primary variables of interest. For film counts, the destination incentive coefficient implies that countries with active incentive schemes receive approximately 18\% more international productions ($p < 0.01$), while the origin coefficient implies that studios reduce overseas filming by about 12\% when their home country introduces incentives ($p < 0.05$).

For budgets, an asymmetry emerges. The destination attraction effect is similar in magnitude (+20\%) but statistically insignificant, while the origin pull-back is much larger ($-38\%$, $p < 0.05$). Home-country incentives appear to most effectively retain high-budget productions, which have the most to gain in absolute terms from a domestic rebate on a large qualifying spend. This pattern---a strong destination effect for film counts but not budgets, and a strong origin effect for budgets but not counts---persists across Models~4--7.

\subsubsection{Incentive Design}

Model~4 (Table~\ref{tab:model4}) disaggregates destination incentives by instrument type. The Olsberg SPI classification groups film incentive schemes into three categories: cash rebates (direct spend-based payments), tax credits (covering both refundable and non-refundable credit structures), and tax shelters (investor-side tax relief mechanisms such as the Belgian and Luxembourg regimes). The specification retains the aggregate destination incentive dummy, with cash rebates as the omitted reference category; the tax credit and tax shelter coefficients therefore capture the differential effect of those instruments relative to a rebate scheme.

For film counts, cash rebates are the most robustly effective instrument: the baseline destination incentive coefficient, which in this specification is identified off rebate-using destinations, is $+16.3\%$ ($p < 0.01$). Tax credits add an insignificant $+8.1\%$ differential (implying a total effect of approximately $+26\%$, but this is not separately significant), and tax shelters show a negative differential of $-9.3\%$ (insignificant), implying a total effect of roughly $+5\%$. The ranking---rebates most effective, tax credits broadly similar, tax shelters weakest---is consistent with the logic that foreign producers need directly-monetisable benefits: a cash rebate pays out regardless of tax liability, while a tax shelter requires specific domestic investor structures that may be difficult for foreign productions to access.

For budgets, the type heterogeneity is far sharper. Cash rebates show a large positive coefficient ($+30\%$) but it is not statistically significant, reflecting the noisier budget-weighted margin. Tax credits add a small negative differential ($-8\%$, insignificant). Tax shelters, however, are associated with a dramatic $-67\%$ differential ($p < 0.01$), implying that destinations offering only tax shelter regimes attract roughly $56\%$ less inward production spending than destinations with no scheme at all. This pattern is consistent with tax shelter regimes being structured around investor tax optimisation rather than genuine production attraction, and potentially selecting smaller, co-production-heavy portfolios. Selection effects---tax-shelter countries (notably Belgium and Luxembourg) being small economies without the infrastructure for large-budget productions---likely contribute and cannot be separated at this specification.

\subsubsection{Generosity}

Model~5 (Table~\ref{tab:model5}) replaces dummies with continuous generosity rates. At the mean active generosity of 27.6\%, the implied effects are $+17\%$ for destination film counts and $-13\%$ for origin film counts, aligning closely with the binary estimates from Model~3 ($+18\%$ and $-12\%$ respectively). The consistency across functional forms strengthens confidence in the core estimates and confirms that the incentive rate matters proportionally: a 30\% credit attracts more production than a 20\% credit.

\subsubsection{Duration and Timing}

Models~6 and~7 (Tables~\ref{tab:model6} and~\ref{tab:model7}) test whether incentive effects accumulate over time. Model~6 finds no evidence of duration effects: each additional year of incentive operation adds approximately 0.6\% to inward film production, but this is statistically insignificant. The benefit of having incentives appears to be a one-time level shift rather than a compounding advantage, providing no support for the ``cluster-building'' narrative.

Model~7, however, reveals that \textit{when} a country adopted incentives matters. Early adopters (pre-2005) attract 26\% more films and 45\% more production spending, both highly significant. Late adopters attract only 15\% more films ($p < 0.05$) and show no significant budget effect (+10\%, $p = 0.53$). This decomposition resolves an apparent puzzle from Model~3: the insignificant aggregate budget destination effect was masking two very different effects that averaged out. Early adopters attracted both more films and substantially more spending; late adopters diluted the pooled estimate.

The distinction between Models~6 and~7 is important. The absence of duration effects means incentive programmes do not become more effective the longer they operate. But early adopters did better---not because their incentives compounded over time, but because they entered a less competitive landscape. As more countries introduced schemes, the marginal attraction of any single destination's incentive was diluted.

% ----------------------------------------------------------------
\subsection{Identification: Country-Pair Fixed Effects}
\label{sec:pair_fe_results}
% ----------------------------------------------------------------

Tables~\ref{tab:model8}--\ref{tab:model11} report pair fixed effects estimates, which absorb all time-invariant bilateral heterogeneity and identify incentive effects exclusively from within-pair temporal variation.

\subsubsection{Pair FE Incentive Effects}

Model~8 (Table~\ref{tab:model8}) provides the most demanding test. For film counts, the destination incentive effect survives at approximately $+9\%$ ($p < 0.01$), attenuated from the PPML estimate of $+18\%$. The attenuation is expected: pair FE strip out the cross-sectional correlation between incentive adoption and pre-existing production relationships. The roughly 50\% reduction from PPML to pair FE quantifies the selection component---countries with established film industries were more likely to adopt incentives, inflating the cross-sectional estimate. The portion that survives ($+9\%$) represents the genuine within-pair effect.

This result is robust across subsamples. Excluding Anglo-to-Anglo pairs yields $+9.4\%$; excluding all Anglo countries yields $+7.2\%$; dropping EFW controls yields $+9.4\%$ (all $p < 0.01$). The destination attraction effect for film counts is not an artefact of the Anglo production network.

For budgets, the pair FE destination effect is larger ($+22\%$, $p < 0.05$) and significant in the full sample and excluding Anglo-to-Anglo pairs ($+23\%$, $p < 0.05$). However, it collapses entirely when all Anglo countries are excluded ($-0.0\%$). The budget-weighted destination effect is driven by Anglo destinations---primarily the UK, Canada, and Australia attracting US production spending.

The origin pull-back for film counts is $-8.8\%$ ($p < 0.001$) in the full sample but vanishes entirely in the excluding-all-Anglo subsample ($-0.8\%$, insignificant). The PPML origin estimate was inflated by cross-sectional differences between Anglo and non-Anglo countries. When Anglo home countries introduced incentives, their studios---which account for a disproportionate share of all international production---pulled activity back. Non-Anglo home-country incentives produce no measurable reduction in outward filming. For budgets, the origin effect is insignificant across all pair FE subsamples, indicating that the large PPML estimate ($-38\%$) reflected between-pair rather than within-pair variation.

\subsubsection{Event Study}

Model~9 (Table~\ref{tab:model9}) replaces the destination incentive dummy with relative-time dummies to examine dynamic effects and test parallel trends.

For film counts, the full-sample pre-treatment coefficients are not perfectly flat (average 0.075, with $\tau_{-5}$ marginally significant), but the pattern is erratic rather than monotonically trending. Post-treatment coefficients show a clear jump at $\tau_{+1}$ and $\tau_{+2}$ ($p < 0.01$), consistent with a one-to-two year lag between incentive introduction and production decisions responding. The incremental effect---the post-treatment average (0.155) minus the pre-treatment average (0.075)---is approximately 0.08, aligning with the Model~8 pair FE estimate (0.091). However, without Anglo countries, the pre-treatment coefficients are larger and individually significant, with the post-treatment average (0.112) falling \textit{below} the pre-treatment average (0.145). This undermines causal claims for the non-Anglo film count subsample.

The budget event study provides the cleanest causal evidence in this paper. The full-sample pre-treatment average is essentially zero (0.007), with no individually significant coefficients. Post-treatment effects are small initially but grow substantially from $\tau_{+5}$ onward, consistent with high-budget productions' longer planning horizons. The excluding-Anglo-to-Anglo subsample is the strongest specification: a slightly negative pre-treatment average ($-0.069$), followed by large and increasingly significant post-treatment effects ($\tau_{+6}$: 0.726, $p < 0.01$; $\tau_{+8}$: 1.067, $p < 0.01$; $\tau_{+9}$: 0.909, $p < 0.01$). The growing trajectory suggests incentives first attract smaller, more mobile productions that appear quickly in film counts, while larger-budget activity follows with a multi-year lag as studios develop projects, scout locations, and observe successful early productions.

\subsubsection{Duration Under Pair FE}

Model~10 (Table~\ref{tab:model10}) confirms the null duration result from Model~6 under pair FE identification. The years-active coefficient is insignificant for film counts across all subsamples (0.1--0.4\% per year, all $p > 0.1$) and for budgets in three of four subsamples. The one exception is the excluding-all-Anglo budget subsample, where years active is significantly \textit{negative} ($-6.9\%$ per year, $p < 0.01$), suggesting that non-Anglo destinations may see incentive effectiveness decay as competitive entry erodes their initial advantage. However, the small sample (835 observations, 72 pairs with variation) warrants caution.

\subsubsection{Incentive Type Heterogeneity Under Pair FE}

Model~11 (Table~\ref{tab:model11}) subjects the incentive type decomposition from Model~4 to pair FE identification, testing which instrument types retain a genuine within-pair effect once cross-sectional selection is removed.

For film counts, cash rebates---the reference category, captured by the destination incentive coefficient---retain a significant positive within-pair effect of $+8.5\%$ in the full sample ($p < 0.01$), attenuated from the PPML estimate of $+16.3\%$ but broadly consistent with the Model~8 aggregate estimate. The effect is stable across subsamples: $+10.0\%$ excluding Anglo-to-Anglo pairs, $+8.1\%$ excluding all Anglo countries, and $+8.6\%$ without EFW controls (all $p < 0.05$). Tax credits add a small positive differential ($+6.9\%$ in the full sample) that is statistically insignificant in every specification, and flips to slightly negative in the non-Anglo subsample ($-5.0\%$, insignificant). This pattern is consistent with tax credits and rebates producing broadly similar within-pair effects once cross-sectional selection is absorbed; the cross-sectional differences visible in Model~4 appear largely attributable to selection. Tax shelters show a consistently negative differential---$-15.8\%$ in the full sample and $-15.4\%$ excluding EFW (both $p < 0.05$)---implying a total tax shelter effect near zero or slightly negative, though the differential loses significance in the non-Anglo subsample.

For budgets, the pair FE results echo the PPML type hierarchy more cleanly than for film counts. Cash rebates show a substantial positive effect in the full sample ($+25.3\%$, $p < 0.05$) and excluding Anglo-to-Anglo pairs ($+27.3\%$, $p < 0.05$), though both estimates collapse to insignificance in the excluding-all-Anglo subsample ($+4.0\%$). Tax credits show essentially zero differential across specifications, indicating that credits and rebates have indistinguishable within-pair budget effects once the base rebate effect is accounted for. Tax shelters, consistent with Model~4, carry a large negative differential---$-63.1\%$ in the full sample, $-57.5\%$ excluding Anglo-to-Anglo, and $-66.6\%$ without EFW (all $p < 0.01$). The magnitude and consistency of this result across specifications, and its survival under pair FE identification, is the strongest instrument-design finding in the paper: within bilateral pairs, destinations operating tax shelter regimes see markedly less inward production spending than destinations with active cash rebate schemes.

These results refine the policy implications of Model~4. Under strict within-pair identification, the cross-sectional differences between cash rebates and tax credits largely disappear---both are similarly effective at attracting international production. The key instrument-design lesson is negative: tax shelter regimes, as implemented in the available data, do not function as effective production attraction mechanisms for either extensive-margin film counts or intensive-margin production spending.

% ----------------------------------------------------------------
\subsection{Cross-Dataset Validation: BaTIS Audiovisual Services}
\label{sec:batis_results}
% ----------------------------------------------------------------

The parallel estimation on OECD--WTO Balanced Trade in Services data for audiovisual services (2005--2023) provides independent validation from a dataset with different units, coverage, and measurement error. The BaTIS results are reported in full in the appendix tables; this section summarises the key convergences and divergences with the IMDb findings.

The gravity framework fits BaTIS more closely than IMDb in several respects. Both origin and destination GDP are large and highly significant (elasticities of approximately 0.9 and 0.7 respectively), consistent with standard gravity predictions for services trade. Distance is larger ($-0.52$ to $-0.70$) and closer to goods trade estimates than the IMDb elasticities ($-0.18$ to $-0.28$). Common language doubles audiovisual services flows ($+102\%$), exceeding even the large IMDb effect ($+69\%$). Colonial ties and RTAs are both positive and significant in BaTIS, reversing the negative IMDb findings---a divergence that likely reflects BaTIS capturing ongoing post-colonial media trade and the benefits of trade liberalisation for broader services flows, rather than discrete production relocation decisions.

The core incentive findings converge across datasets. The destination incentive effect under pair FE is $+27\%$ in BaTIS ($p < 0.01$), larger than the IMDb film count effect ($+9\%$) but qualitatively consistent. Crucially, the BaTIS destination effect is stable across Anglo subsamples ($+26\%$ excluding all Anglo countries), providing stronger evidence than IMDb that the destination attraction effect generalises beyond the Anglo production network. The origin pull-back observed in IMDb PPML ($-12\%$) is also present in BaTIS PPML ($-15\%$), though under pair FE the BaTIS origin effect flips to \textit{positive} ($+22\%$, $p < 0.01$). This divergence is interpretable: home-country incentives may simultaneously pull film shoots back (reducing IMDb flows) while deepening broader audiovisual services relationships (increasing BaTIS flows) through expanded post-production, distribution, and co-financing activity.

BaTIS provides mixed evidence on duration effects. In PPML, years active is positive and significant ($+2.2\%$ per year, $p < 0.05$), unlike the null IMDb result, suggesting that official services trade captures infrastructure accumulation that discrete film counts do not. However, under pair FE the coefficient reverses to significantly \textit{negative} ($-2.1\%$ per year, $p < 0.01$), indicating that the PPML duration effect reflected cross-sectional selection rather than genuine within-pair accumulation. The BaTIS event study is uninformative due to near-perfect multicollinearity among event-time indicators in the balanced panel structure \citep{borusyak2024}.

The incentive type results under pair FE (Model~11) show both striking convergences and an informative divergence between the two datasets. Cash rebates---identified by the destination incentive coefficient in this specification---show a large positive within-pair effect in BaTIS ($+27\%$, $p < 0.01$, full sample) that is stable across all subsamples including the excluding-all-Anglo specification ($+26\%$, $p < 0.01$). This magnitude substantially exceeds the IMDb film count estimate ($+9\%$) but is directionally consistent and, crucially, robust to Anglo-country exclusion in both datasets. Tax credit differentials are negative and significant in BaTIS ($-25\%$ to $-44\%$, $p < 0.01$ across subsamples), implying that the total within-pair effect of a tax credit scheme on BaTIS audiovisual services exports is approximately zero in the full sample and meaningfully negative when Anglo countries are excluded---a sharp divergence from the small positive or null IMDb film count effect for tax credits.

The tax shelter finding is the most striking divergence between datasets. In IMDb, tax shelter destinations attract fewer films and substantially less production spending under pair FE. In BaTIS, tax shelter differentials are \textit{positive} and highly significant across all subsamples ($+34\%$, $p < 0.01$), implying a total within-pair effect for tax shelters of approximately $+61\%$ on audiovisual services exports. The divergence is interpretable if the two dependent variables capture different margins: IMDb filming location data measures where physical production occurs, while BaTIS captures the full bilateral audiovisual services flow including post-production, rights trading, and co-financing. Tax shelter regimes, which are typically structured around investor tax relief tied to financing rather than physical production, may channel substantial bilateral services trade through specific financial structures without producing a proportionate increase in observable filming activity. This reading is consistent with the Belgian and Luxembourg tax shelter systems in particular, which are widely used as co-financing vehicles.

Taken together, the cross-dataset pattern refines the policy message: cash rebates are the most broadly effective instrument for attracting observable inward production across both datasets and all subsamples; tax credits perform comparably to rebates in IMDb but weaker than rebates in BaTIS services trade; and tax shelters operate as financing-channel instruments rather than production-attraction instruments, boosting services flows without a commensurate rise in physical production activity.

% ----------------------------------------------------------------
\subsection{Interpretation}
\label{sec:interpretation}
% ----------------------------------------------------------------

Five themes emerge from the results that warrant explicit discussion, because the economic interpretation is not always obvious from the coefficients alone.

\subsubsection{Selection, Not Causation, Explains Much of the Incentive Effect}

The single most important finding is quantitative rather than qualitative: roughly half of the headline destination incentive effect is selection, not treatment. The PPML estimate of $+18\%$ more international productions in IMDb data (Model~3) falls to $+9\%$ once pair fixed effects strip out the time-invariant bilateral heterogeneity (Model~8). The BaTIS data show a parallel collapse in the origin effect, which even flips sign under pair FE (from $-15\%$ to $+22\%$). These attenuations are exactly what one expects if countries with established film industries were disproportionately likely to adopt incentive schemes. A country that already hosts substantial inbound production has both the political constituency and the administrative infrastructure to enact incentives; a country without either will find it harder to legislate a rebate, harder to administer it, and harder to generate the supply-side response that would make it effective. The cross-sectional correlation between incentive presence and production activity therefore contains a large component that would exist regardless of whether the incentive itself caused anything.

This is an important cautionary reading for the existing film-incentive literature, which has relied heavily on cross-sectional or state-panel variation \citep{button2019, thom2018}. Without absorbing bilateral pair heterogeneity, the estimated incentive effect will tend to reflect both a genuine treatment effect and the institutional capacity that makes adoption likely in the first place. The within-pair effect that survives ($+9\%$ for film counts, $+27\%$ for audiovisual services) is still economically meaningful, but is substantially smaller than the numbers typically reported in policy briefs and industry advocacy documents.

\subsubsection{The Anglo Production Network Is a Pervasive Confound}

Several results change materially when the five major Anglo countries (US, UK, Canada, Australia, New Zealand) are excluded from the sample. The origin pull-back for IMDb film counts falls from $-9\%$ (significant) to $-0.8\%$ (insignificant). The pair FE destination incentive effect for IMDb budgets collapses from $+25\%$ in the full sample to an insignificant $+4\%$ when all Anglo countries are excluded. The event study's post-treatment average for film counts falls below the pre-treatment average, undermining the causal interpretation for the non-Anglo sample. And in Model~11, the tax shelter differential on IMDb film counts loses significance once Anglo pairs are excluded, though the IMDb film count effect for cash rebates and the BaTIS destination effect are both stable across Anglo subsamples.

The explanation is likely structural rather than statistical. Hollywood studios produce the overwhelming majority of internationally mobile productions, and the bilateral flows between Hollywood and the English-speaking commonwealth countries (especially UK and Canada) account for a disproportionate share of observed international filming in any given year. When a major Anglo production hub like Canada or the UK introduces a rebate or credit, the near-term response is dominated by US studios redirecting projects into that hub---a genuine causal effect, but one that reflects the existing concentration of mobile production in a narrow network rather than a general response of ``international producers'' to ``destination incentives.'' The non-Anglo subsample is effectively asking a different question: what happens when, say, Hungary or South Korea introduces a rebate and attracts production from Germany, France, or Japan? The answer appears to be that the effect is positive but smaller, and noisier for the budget margin in particular, reflecting both smaller baseline flows and a less concentrated producer population that is not monitoring and responding to subsidy competition in the same way.

This has a methodological implication worth emphasising: aggregate estimates of incentive effectiveness that pool Anglo and non-Anglo relationships risk projecting the dynamics of the Hollywood-Commonwealth corridor onto the global film trade, where they do not generalise cleanly. Policy recommendations derived from such aggregates may not transfer to non-Anglo destinations contemplating similar schemes.

\subsubsection{The Two Datasets Measure Different Things---and This Is Informative}

The sharp divergence between IMDb and BaTIS on tax shelter effects---strongly negative in IMDb, strongly positive in BaTIS---is not a contradiction but a reminder that ``film production'' is a bundle of activities, and different data sources capture different slices of that bundle. IMDb filming location data is the best available proxy for where cameras actually rolled. BaTIS captures the full value of bilateral audiovisual services exports: not only physical production but also post-production, visual effects, rights trading, distribution, co-financing, and licensing flows. Tax shelter regimes, as typically designed (the Belgian and Luxembourg systems are the clearest cases), are structured to channel investor capital into audiovisual productions in exchange for tax relief; they do not necessarily require physical production to occur on the shelter country's territory. A Belgian tax shelter investment in a French production generates bilateral services flows without generating a Belgian filming location.

Read this way, the tax shelter finding is a feature, not a bug, of the cross-dataset comparison. It demonstrates that incentive instruments target different margins of the audiovisual value chain, and the appropriate evaluation metric depends on what the policy is actually trying to achieve. If the goal is attracting physical filming (and the associated local crew employment, location spending, hospitality, and so on), cash rebates are the demonstrated winner and tax shelters are ineffective or worse. If the goal is increasing the country's presence in the broader audiovisual services economy through co-production and financing relationships, tax shelters appear to do exactly that. Conflating these two objectives---as much of the policy discussion does---leads to imprecise evaluation.

A related observation applies to the origin pull-back divergence. The IMDb origin effect is negative (studios film less abroad when their home country introduces incentives), while the BaTIS origin effect is positive under pair FE (home-country incentives are associated with \textit{more} bilateral audiovisual services trade). Both can be true simultaneously. Physical filming may be pulled back to the home country, while services-trade linkages with partner countries---post-production contracts, co-financing arrangements, licensing deals---deepen precisely because the home-country scheme provides the fiscal foundation for larger productions that then trade services bilaterally. The two effects operate on different margins of the value chain and should not be expected to move together.

\subsubsection{Why Film Gravity Looks Different from Goods Gravity}

The gravity controls in Models~1 and~2 deserve interpretation beyond the incentive focus. Distance elasticities are substantially smaller than in goods trade ($-0.24$ to $-0.28$ for film counts, compared to $-0.7$ to $-1.5$ typical for goods). Common language effects are substantially larger ($+69\%$ for film counts, $+102\%$ for BaTIS services trade, versus $+15\%$ to $+30\%$ in goods). Contiguity is positive and significant but modest. Colonial ties and regional trade agreements are \textit{negative} for IMDb film production, reversing the standard gravity prediction.

These patterns are internally coherent if one takes seriously that film is a language-intensive service, not a physical commodity. Shipping a physical good over long distances is expensive; shipping a production crew is also expensive, but the fixed costs of mounting a production dwarf the marginal transport cost, so distance matters less. Filming in a country where the language is unfamiliar imposes persistent, production-spanning costs---translation, directorial miscommunication, post-synchronisation---that a twenty-hour flight simply does not. The large language coefficient captures this friction directly.

The negative colonial coefficient in IMDb is the most surprising result and probably reflects two coincident effects. First, post-colonial cultural production in many ex-colonies has intentionally diverged from the metropole's film tradition (Indian cinema, Latin American cinema, the Nigerian industry) in ways that weaken rather than strengthen bilateral production relationships with former colonisers. Second, colonial-era relationships are substantially overlapped with Anglo production dynamics: the excluded colonial pairs in the non-Anglo subsamples absorb much of this effect. The positive colonial coefficient in BaTIS services trade, by contrast, likely captures ongoing post-colonial media trade (film distribution, broadcast rights, licensed content) that survives the decoupling of physical production.

\subsubsection{Timing Trumps Duration}

Models~6 and~7 together tell a story that is worth stating explicitly because it contradicts a standard industry narrative. The ``cluster-building'' argument holds that incentives need time to work: each additional year builds infrastructure, trains crews, and attracts complementary businesses, generating compounding effects. This argument is not supported by the data. Years of active operation add approximately zero to inward production in both IMDb and BaTIS under pair FE identification. The BaTIS result is actually negative ($-2\%$ per year under pair FE), consistent with competitive entry eroding first-mover advantage.

What does matter is \textit{when} a country adopted. Early adopters (pre-2005) attract substantially more production than late adopters, for both film counts and budgets. The difference is not that their incentives became more effective over time but that they entered an uncrowded subsidy market. As more countries introduced schemes, the marginal subsidy-competitive advantage of any single destination eroded. This is a first-mover effect in a global subsidy competition rather than a cluster-accumulation effect.

The policy implication is uncomfortable for late adopters. A country introducing incentives today is not facing the same effectiveness distribution that early adopters faced; it is entering a market in which incentives have become a minimum entry requirement rather than a distinguishing advantage. The absence of duration effects means that sustaining a scheme does not generate a compounding return; the first-mover advantage accrued to those who entered before subsidy competition became pervasive, and it cannot be recovered by a latecomer simply by maintaining a comparable scheme for long enough.

\subsubsection{Synthesis}

The most defensible overall claim the evidence supports is modest: film production incentives do cause some bilateral production to occur that would not otherwise occur, but the causal effect is roughly half the size of cross-sectional estimates, is concentrated in the Anglo production network, depends substantially on instrument design (cash rebates outperforming tax shelters for physical production), and does not compound with time. Much of what was previously read as an incentive effect in the literature is, on this evidence, a selection effect reflecting which countries were positioned to adopt incentives in the first place. The cross-dataset divergence on tax shelters further suggests that ``film incentive effectiveness'' is not a single parameter but a set of distinct margins---physical production, services trade, co-financing---that different instruments target differently. Aggregate evaluations that do not distinguish these margins, or that do not absorb bilateral selection, are likely to both overstate the overall effectiveness of incentives and misattribute effects across instrument types.


